% !PRINTED 286
\weberpage{286}

%<<CONT
geschränkt der geometrischen Anschauung und der Bezeichnungsweise der Differentialrechnung.

Wir betrachten zunächst zwei reelle Functionen $\varphi(x,y)$, $\psi(x,y)$, und deuten die reellen Variablen $x, y$ als rechtwinklige Coordinaten in einer Ebene. Die Gleichungen
\weq{1}{\varphi(x,y) = 0,\qquad \psi(x,y) = 0}
stellen dann zwei Curven dar, die wir kurz die Curve $\varphi$ und die Curve $\psi$ nennen. Wir nehmen zunächst an, diese Curven seien geschlossen und erstrecken sich nicht ins Unendliche. Wir nehmen ausserdem an, $\varphi$ sei im \gesperrt{Inneren} der Curve $\varphi$ \gesperrt{negativ}, im \gesperrt{Aeusseren} \gesperrt{positiv}; ebenso $\psi$ im Inneren von $\psi$ negativ, im Aeusseren positiv.

Bezeichnen wir, wie die Fig. 6 zeigt, den Winkel, den die Normale $n$ in irgend einem Punkte von $\varphi$, in der Richtung, in der $\varphi$ wächst, also nach aussen gezogen, mit der Richtung der positiven $x$-Axe bildet, mit $\vartheta$, so ist, wenn $\varphi'(x)$ und $\varphi'(y)$ die partiellen Ableitungen von $\varphi$ sind, und die Quadratwurzel positiv genommen ist,
\[
\cos\vartheta = \frac{\varphi'(x)}{\sqrt{\varphi'(x)^2 + \varphi'(y)^2}},\qquad \sin\vartheta = \frac{\varphi'(y)}{\sqrt{\varphi'(x)^2 + \varphi'(y)^2}}.
\]
Ziehen wir die Tangente $t$ so, dass sie zu der eben bezeichneten Normale $n$ so liegt, wie die positive $y$-Axe zur positiven $x$-Axe und bezeichnen den Winkel, den sie mit der positiven $x$-Axe bildet, mit $\xi$, so ist
\[
\cos\xi = -\frac{\varphi'(y)}{\sqrt{\varphi'(x)^2 + \varphi'(y)^2}},\qquad \sin\xi = \frac{\varphi'(x)}{\sqrt{\varphi'(x)^2 + \varphi'(y)^2}}.
\]
Bezeichnen wir den Fortschritt auf $\varphi$ in der Richtung $t$ als den \gesperrt{positiven} und sind $dx, dy$ die Projectionen dieses Fortschrittes, so sind $dx$ und $dy$ proportional und im Zeichen übereinstimmend mit $-\varphi'(y)$, $\varphi'(x)$, und wenn $\Phi$ eine beliebige andere Function ist, so ist
\[
d\Phi = \Phi'(x)\,dx + \Phi'(y)\,dy
\]
im Vorzeichen übereinstimmend mit der \gesperrt{Functionaldeterminante}
\weq{2}{[\varphi, \Phi] = \varphi'(x)\Phi'(y) - \varphi'(y)\Phi'(x).}
Bei der üblichen Annahme über die Coordinatenrichtung ist die positive Fortschrittsrichtung die, bei der das Innere der Fläche zur Linken liegt.
